\documentclass[Chp3.tex]{subfiles}

\begin{document}


\section{Ranking wealth distributions using choice under ambiguity}
\label{sec:ExtendedModel}

\par Consider the excerpt from \cite{g09} about the choice under ambiguity model, which motivates the modeling choices that will be made to characterize a ranking over wealth distributions\footnote{I move away from the notation used there in favor of the \say{choice over lotteries} notation to tie into the birth lottery analogy made earlier in the paper}:

\begin{quote}


To make sure that we understand the structure, observe that there are two sources of uncertainty: the choice of the state $s$, which is sometimes referred to as \say{subjective uncertainty}, because no objective probabilities are given on it, and the choice of $y$, which is done with objective probabilities once you chose your act and Nature chose a state. Specifically, if you choose $p \in L$ and Nature chooses $s \in S$, a roulette wheel is spun, with distribution $p_s$ over the outcomes $Y$, so that your probability to get outcome $y$ is $p_s(y)$.

\end{quote}

\subsection{Analytical framework}

\par Denote the set of outcomes $Y = [0,\bar{y}]$. Then, the choice set is given by:

$$ L = \bigg\{ p: 2^{[0,\bar{y}]} \to \mathbb{R} \hspace{2mm} | \hspace{2mm} p(\cdot) \hspace{2mm} \text{is an income frequency distribution}  \bigg\}. $$

\par We may define a binary relation over both sets $Y,L:$

$$ \succsim_{y} \subseteq [0,\bar{y}] \times [0,\bar{y}] \subseteq \mathbb{R} \times \mathbb{R} $$

$$ \succsim \; \subseteq L \times L. $$

\par Notice that $\succsim_y$ may be represented by a real-valued utility function. This is the social welfare $U(y)$ which is a key object of analysis in this paper. 

\subsubsection*{Preliminary remarks}

We are concerned with the comparison of two frequency distributions $f(w)$ of an outcome $w$ which we refer to as wealth. We seek to use the notion of \textit{uncertainty aversion} as an analogy to the use of the notion of risk aversion in the characterization of a ranking over income distributions.

The presence of both objective and subjective uncertainty is at the heart of this analysis. Section 2 covered the analysis for objective uncertainty. Thus, the wealth frequency distribution $f(w)$ must be formalized in this abstract setting so that it explicitly captures both forms of uncertainty. Namely, each wealth distribution is associated with some relevant, underlying state space $S$ (the source of subjective uncertainty), and its objective component $p(y) \in L$. 

\par This suggests that the objects of interest (rankable wealth distributions $f(w)$) are a family of distributions $p(y) \in L$, indexed by elements of the state space $s \in S$, or:

$$\{ p_s(y) \}_{s \in S}.$$

\par Thus, we may view each act $f \in F$ as a \say{state-contingent frequency distribution}. Moreover, upon fixing a state $s' \in S$, the analysis will become identical to the use of the choice under objective uncertainty model that was used to reach a ranking over income frequency distributions.

With this in mind, we work under the following assumption on the functional form of wealth distributions for the remainder of this paper:

\begin{assu}
Each wealth distribution $f(w)$ can be written as $p_s(y)$.
\end{assu}


\subsection{Axiomatization}

\begin{ax.AA}[Weak Order]
$\succsim$ is complete and transitive.
\end{ax.AA}

\begin{ax.AA}[Continuity]
For every $f,g,h \in F$,  if $f \succ g \succ h $, there exists $\alpha, \beta \in (0,1)$ such that 
$$ \alpha f + (1-\alpha) h \succ g \succ \beta f + (1-\beta) h .$$ 
\end{ax.AA}

\begin{ax.C}[C-Independence]
For every $f,g \in F$, every constant $h \in F$ and every $\alpha \in (0,1)$,
$$ f \succsim g \iff \alpha f + (1-\alpha)h \succsim \alpha g + (1-\alpha)h. $$
\end{ax.C}

\begin{ax.AA}[Monotonicity]
For every $f,g \in F$, $f(s) \geq g(s)$ for all $s \in S$ implies $f \geq g$.
\end{ax.AA}

\begin{ax.AA}[Non-trivality]
There exists $f,g \in F$ such that $f \succ g$.
\end{ax.AA}

\begin{ax.U}[Uncertainty Aversion]
For every $f,g \in F$, if $f \sim g$, then, for every $\alpha \in (0,1)$,
$$ \alpha f + (1-\alpha)g \succsim f. $$
\end{ax.U}

\subsection{Expected utility representation of the ranking over wealth distributions}

\par  Finally, we have the representation theorem by \cite{gs89}.

\begin{tm}
$\succsim$ satisfies AA1, AA2, C-Independence, AA4, AA5, and Uncertainty Aversion if and only if there exists a closed and convex set of probabilities on $S$, $C \subseteq \Delta(S)$, and a non-constant function $U: Y \to \mathbb{R}$ such that, for every $f, f^* \in F$,
\[
f \succsim f^*
\iff
\min_{\lambda \in C} \int_{S} \big(\mathbb{E}_{p_s}[U]\big)\, d\lambda
\ge
\min_{\lambda \in C} \int_{S} \big(\mathbb{E}_{p_s^*}[U]\big)\, d\lambda.
\]
\end{tm}


From here on out, we assume that wealth distributions will be ranked according to:

$$ W^{A} \equiv \min_{\lambda \in C} \int_{S} (\mathbb{E}_{p_s}[U]) d \lambda $$

$$ = \min_{\lambda \in C} \int_{S} \int_{0}^{\bar{y}} p_s(y) U(y) dy d\lambda. $$

\par \textit{Notation.} Superscript $A$ denotes the ambiguity specification; this is the first appearance of $W^{A}$ in the main text.




\end{document}
